\documentclass[../../main.tex]{subfiles}
\begin{document}

\begin{flushright}
\footnotesize\textit{【自動文字起こし・要確認】}
\end{flushright}

\noindent
{\bf [解]}
\paragraph{(1)}

$f(x)=\dfrac{1}{2}x\{1+e^{-2(x-1)}\}$である．
\begin{align}
f'(x)&=\dfrac{1}{2}\left[(1+e^{-2(x-1)})+x(-2)e^{-2(x-1)}\right]\\
&=\dfrac{1}{2}\left[1+(1-2x)e^{-2(x-1)}\right]
\end{align}
である．$\dfrac{1}{2}<x$の時，$(1-2x)e^{-2(x-1)}<0$だから，$f'(x)<\dfrac{1}{2}\quad\cdots①$である．
\begin{align}
f''(x)&=\dfrac{1}{2}e^{-2(x-1)}\left[-2+(1-2x)(-2)\right]\\
&=-e^{-2(x-1)}(2-2x)
\end{align}
から下表を得る．
\begin{center}
\begin{tabular}{c|ccccc}
$x$ & $\dfrac{1}{2}$ & $\cdots$ & $1$ & $\cdots$ & $\infty$\\
\hline
$f''$ & & $-$ & $0$ & $+$ & \\
$f'$ & & $\searrow$ & $0$ & $\nearrow$ & \\
\end{tabular}
\end{center}
したがって，$\dfrac{1}{2}<x$の時，$0\le f'(x)\quad\cdots②$である．①②から示された．$\blacksquare$

\paragraph{(2)}

$f(1)=1$だから，
\begin{align}
x_{n+1}-1=f(x_n)-f(1)\quad\cdots③
\end{align}

$x_n=1$なる$n$がある時，③から$n\le m$なる任意の$m$に対し，$f(x_m)=1$だから，$x_n\to1\ (n\to\infty)$である．

以下，すべての$n$で$x_n\ne1$の場合を考える．$(1)$の表から$x>\dfrac{1}{2}$の時$f(x)$は単調増加で$f\left(\dfrac{1}{2}\right)=\dfrac{1}{4}(1+e)>\dfrac{1}{2}$だから，$x>\dfrac{1}{2}\implies f(x)>\dfrac{1}{2}$である．$x_0>\dfrac{1}{2}$だから，帰納的に任意の$n$に対し$x_n>\dfrac{1}{2}$である．

このことに注意すると，平均値の定理から（$\because f(x)$が微分可能である），
\begin{align}
f(x_n)-f(1)=f'(c)(x_n-1)\quad\left(c\text{は}\dfrac{1}{2}<c\text{，}1\text{と}x_n\text{の間}\right)
\end{align}
なる$c$がある．③から，
\begin{align}
|x_{n+1}-1|=|f'(c)||x_n-1|<\dfrac{1}{2}|x_n-1|\quad(\because(1))
\end{align}

くり返し用いて，
\begin{align}
|x_n-1|<\left(\dfrac{1}{2}\right)^n|x_0-1|\longrightarrow0\quad(n\to\infty)
\end{align}

はさみうちから
\begin{align}
x_n\longrightarrow1\quad(n\to\infty)
\end{align}

\bigskip

\end{document}
